\chapter{Fuzzy Distributions and the Linear Duality of a Fuzzy Cosmos}
\label{ch:fuzzy-distributions-linear-duality}

\section{Introduction}
\label{sec:fuzzy-distributions-introduction}

The purpose of this chapter is to construct a fuzzy-domain-theoretic analogue
of the linear double-dual distribution monad.  The construction is internal to
the symmetric monoidal closed category of fuzzy domains
\[
    \mathbf{FDom}_{\mathbb V}.
\]
Thus the relevant duality is not the raw presheaf duality
\[
    [X^{\op},\underline{\mathbb V}],
\]
nor the raw observation object
\[
    [X,\underline{\mathbb V}].
\]
Instead, the duality is the internal hom duality of fuzzy domains.

The dualizing object is the self-enriched truth-value object
\[
    \underline{\mathbb V},
\]
regarded as a fuzzy domain in the universe-bounded sense used below.  For a
fuzzy domain \(X\), its object of observations is the fuzzy Scott topology
\[
    \mathcal O_{\mathbb V}(X)
    =
    [X,\underline{\mathbb V}]_{\mathrm{ff}},
\]
where the internal hom is taken in
\[
    \mathbf{FDom}_{\mathbb V}.
\]
Thus the objects of
\[
    \mathcal O_{\mathbb V}(X)
\]
are precisely the truth-valued fuzzy-domain morphisms
\[
    X\to\underline{\mathbb V}.
\]
Equivalently, they are the fuzzy Scott opens of \(X\): upper fuzzy predicates
which preserve fuzzy ideal joins.

The fuzzy distribution object of \(X\) is then the internal linear dual of its
fuzzy Scott topology:
\[
    \mathbb D_{\mathbb V}X
    =
    [\mathcal O_{\mathbb V}(X),\underline{\mathbb V}]_{\mathrm{ff}}.
\]
Equivalently,
\[
    \mathbb D_{\mathbb V}X
    =
    [[X,\underline{\mathbb V}]_{\mathrm{ff}},
      \underline{\mathbb V}]_{\mathrm{ff}}.
\]
This is the fuzzy Scott double-dual distribution construction.

The object assignment is therefore
\[
    X
    \longmapsto
    X^{\perp}
    =
    [X,\underline{\mathbb V}]_{\mathrm{ff}}
    \longmapsto
    X^{\perp\perp}
    =
    [[X,\underline{\mathbb V}]_{\mathrm{ff}},
      \underline{\mathbb V}]_{\mathrm{ff}}.
\]
The first assignment is contravariant, and applying it twice gives the
covariant endofunctor
\[
    \mathbb D_{\mathbb V}
    =
    (-)^{\perp\perp}.
\]

\paragraph{Standing size convention.}
We work relative to the same tower of universes used in the preceding chapters.
When
\[
    \underline{\mathbb V}
\]
or an internal hom
\[
    [X,\underline{\mathbb V}]_{\mathrm{ff}}
\]
is not small in the base universe, it is interpreted in the next universe of
the fixed tower.  Formally,
\[
    \mathbf{FDom}_{\mathbb V}^{+}
\]
denotes the full category of fuzzy domains small in the next universe, with the
same notion of finite-flat-colimit-preserving morphism.  It is used only to
ensure that the internal homs appearing in this chapter exist as objects of the
ambient universe.

\paragraph{Standing closed structure.}
We use the symmetric monoidal closed structure
\[
    A\boxtimes_{\mathrm{ff}}B
    \dashv
    [B,-]_{\mathrm{ff}}
\]
on fuzzy domains.  Thus, for fuzzy domains \(A,B,C\), there is a natural
bijection
\[
    \mathbf{FDom}_{\mathbb V}
    (A\boxtimes_{\mathrm{ff}}B,C)
    \cong
    \mathbf{FDom}_{\mathbb V}
    (A,[B,C]_{\mathrm{ff}}).
\]
The evaluation morphism is
\[
    [B,C]_{\mathrm{ff}}\boxtimes_{\mathrm{ff}}B
    \longrightarrow
    C.
\]

\section{The truth-value object as a fuzzy domain}
\label{sec:truth-value-object-as-fuzzy-domain}

\subsection{The enriched construction}
\label{subsec:truth-value-domain-enriched}

Let
\[
    (\mathbb V,\tensor,I,[-,-])
\]
be a fuzzy cosmos.  Its self-enriched truth-value object is the
\(\mathbb V\)-category
\[
    \underline{\mathbb V}
\]
whose objects are the objects of \(\mathbb V\), and whose hom-objects are
\[
    \underline{\mathbb V}(A,B)=[A,B].
\]

\begin{proposition}[The truth-value object is finite-flat cocomplete]
\label{prop:truth-value-object-fuzzy-domain}
The self-enriched truth-value object
\[
    \underline{\mathbb V}
\]
admits all small weighted colimits whose weights are small.  In particular, it
admits all finite-flat weighted colimits and therefore determines a fuzzy domain
in the universe-bounded sense.
\end{proposition}

\begin{proof}
Let
\[
    W:K^{\op}\to\underline{\mathbb V}
\]
be a small weight, and let
\[
    D:K\to\underline{\mathbb V}
\]
be a small diagram.  Since \(\mathbb V\) is cocomplete and tensoring preserves
small colimits in each variable, the enriched weighted colimit is computed in
\(\mathbb V\) by the coend
\[
    W*D
    =
    \int^{k\in K} W(k)\tensor D(k).
\]
This object satisfies the enriched weighted-colimit universal property in the
self-enrichment of \(\mathbb V\).  Hence all small weighted colimits exist in
\(\underline{\mathbb V}\), and in particular all finite-flat weighted colimits
exist.

The corresponding strict algebra structure for the finite-flat ideal monad is
the map
\[
    \alpha_{\underline{\mathbb V}}:
    \operatorname{Idl}_{\mathbb V}(\underline{\mathbb V})
    \to
    \underline{\mathbb V},
    \qquad
    I\longmapsto I*1_{\underline{\mathbb V}}.
\]
Thus
\[
    \underline{\mathbb V}
\]
is a fuzzy domain in
\[
    \mathbf{FDom}_{\mathbb V}^{+}.
\]
\end{proof}

The affine condition
\[
    I\cong 1_{\mathbb V}
\]
supplies canonical weakening maps
\[
    A\to I
\]
for every \(A\in\mathbb V\).  There is no canonical contraction map
\[
    A\to A\tensor A
\]
unless the monoidal structure is cartesian or otherwise equipped with
diagonals.  Thus the linear duality of this chapter is affine-linear rather
than cartesian.

\begin{definition}[Truth-value fuzzy domain]
\label{def:truth-value-fuzzy-domain}
The \emph{truth-value fuzzy domain} is the self-enriched fuzzy domain
\[
    \Omega_{\mathbb V}
    =
    \underline{\mathbb V}.
\]
When no confusion is possible, we continue to write
\[
    \underline{\mathbb V}
\]
for this object.
\end{definition}

\subsection{The posetal reflection}
\label{subsec:truth-value-domain-posetal}

Let
\[
    \Lambda=(L,\leq,\meet,\join,\tensor,\multimap,\bot,\top)
\]
be a fuzzy truth-value object.  The corresponding truth-value fuzzy domain is
the separated fuzzy preorder
\[
    \Lambda
\]
whose underlying set is \(L\), and whose hom-values are
\[
    \Lambda(a,b)=a\multimap b.
\]

The fuzzy-domain structure is given by the existing weighted joins in \(L\).
In particular, if
\[
    I:\Lambda^{\op}\to\Lambda
\]
is a fuzzy ideal, its join is computed by the corresponding weighted join in
the truth-value preorder.

\section{Fuzzy Scott observations}
\label{sec:fuzzy-scott-observations}

\subsection{The enriched construction}
\label{subsec:fuzzy-scott-observations-enriched}

\begin{definition}[Fuzzy Scott observation object]
\label{def:fuzzy-scott-observation-object}
For a fuzzy domain \(X\), its \emph{fuzzy Scott observation object} is
\[
    \mathcal O_{\mathbb V}(X)
    =
    [X,\underline{\mathbb V}]_{\mathrm{ff}}.
\]
Its objects are fuzzy-domain morphisms
\[
    X\to\underline{\mathbb V}.
\]
We call these \emph{fuzzy Scott observations}, or \emph{fuzzy Scott opens}.
\end{definition}

Thus \(\mathcal O_{\mathbb V}(X)\) is not the raw enriched functor category
\[
    [X,\underline{\mathbb V}],
\]
but the internal hom in fuzzy domains.  Equivalently, it is the full
sub-\(\mathbb V\)-category of
\[
    [X,\underline{\mathbb V}]
\]
spanned by the finite-flat-colimit-preserving functors
\[
    X\to\underline{\mathbb V}.
\]

This agrees objectwise with the fuzzy Scott opens of the preceding development.
Namely, a \(\mathbb V\)-functor
\[
    U:X\to\underline{\mathbb V}
\]
lies in
\[
    [X,\underline{\mathbb V}]_{\mathrm{ff}}
\]
exactly when it preserves finite-flat weighted colimits, equivalently when
\[
    U(\alpha_X(I))
    \cong
    \int^{x\in X} I(x)\tensor U(x)
\]
for every fuzzy ideal
\[
    I:X^{\op}\to\underline{\mathbb V}.
\]
The present chapter uses this Scott-open object as a fuzzy domain, via the
internal hom structure, rather than only as a fuzzy frame.

\begin{proposition}[Fuzzy Scott observations preserve fuzzy ideal joins]
\label{prop:fuzzy-scott-observations-preserve-ideal-joins}
Let \(X\) be a fuzzy domain.  A \(\mathbb V\)-functor
\[
    U:X\to\underline{\mathbb V}
\]
is an object of
\[
    \mathcal O_{\mathbb V}(X)
\]
if and only if it preserves finite-flat weighted colimits, equivalently fuzzy
ideal joins.
\end{proposition}

\begin{proof}
By definition,
\[
    [X,\underline{\mathbb V}]_{\mathrm{ff}}
\]
is the internal hom of fuzzy domains.  Its objects are precisely the
fuzzy-domain morphisms
\[
    X\to\underline{\mathbb V}.
\]
A fuzzy-domain morphism is a \(\mathbb V\)-functor preserving finite-flat
weighted colimits.  Since finite-flat weights are the fuzzy ideals of the
fuzzy-domain doctrine, this is exactly preservation of fuzzy ideal joins.
\end{proof}

\begin{definition}[Pullback of fuzzy Scott observations]
\label{def:pullback-fuzzy-scott-observations}
Let
\[
    f:X\to Y
\]
be a fuzzy-domain morphism.  Pullback of fuzzy Scott observations along \(f\)
is the morphism
\[
    f^\ast:
    \mathcal O_{\mathbb V}(Y)\to\mathcal O_{\mathbb V}(X)
\]
defined by
\[
    f^\ast(V)=Vf.
\]
\end{definition}

\begin{proposition}[Functoriality of fuzzy Scott observations]
\label{prop:functoriality-fuzzy-scott-observations}
The assignment
\[
    X\longmapsto\mathcal O_{\mathbb V}(X)
\]
defines a contravariant functor
\[
    \mathcal O_{\mathbb V}:
    (\mathbf{FDom}_{\mathbb V}^{+})^{\op}
    \longrightarrow
    \mathbf{FDom}_{\mathbb V}^{+}.
\]
\end{proposition}

\begin{proof}
If
\[
    f:X\to Y
\]
is a fuzzy-domain morphism and
\[
    V:Y\to\underline{\mathbb V}
\]
preserves finite-flat colimits, then the composite
\[
    Vf:X\to\underline{\mathbb V}
\]
also preserves finite-flat colimits.  Thus \(f^\ast\) is well-defined on
objects.

The morphism
\[
    f^\ast:
    [Y,\underline{\mathbb V}]_{\mathrm{ff}}
    \to
    [X,\underline{\mathbb V}]_{\mathrm{ff}}
\]
is induced by precomposition.  Since finite-flat colimits in internal homs are
computed pointwise, precomposition preserves finite-flat colimits.  Hence
\(f^\ast\) is a fuzzy-domain morphism.

For composable fuzzy-domain morphisms
\[
    X\xrightarrow{f}Y\xrightarrow{g}Z
\]
one has
\[
    (gf)^\ast=f^\ast g^\ast,
\]
and
\[
    (\id_X)^\ast=\id_{\mathcal O_{\mathbb V}(X)}.
\]
Therefore \(\mathcal O_{\mathbb V}\) is contravariantly functorial.
\end{proof}

\subsection{The posetal reflection}
\label{subsec:fuzzy-scott-observations-posetal}

Let \(X\) be a fuzzy domain over a fuzzy truth-value object \(\Lambda\), with
structure map
\[
    \alpha_X:\operatorname{Idl}_{\Lambda}(X)\to X.
\]
A fuzzy Scott observation is an upper fuzzy predicate
\[
    U:X\to L
\]
satisfying
\[
    U(\alpha_X(I))
    =
    \bigjoin_{x\in X}
    I(x)\tensor U(x)
\]
for every fuzzy ideal
\[
    I:X^{\op}\to\Lambda.
\]
Thus
\[
    \mathcal O_{\Lambda}(X)
    =
    \operatorname{Scott}_{\Lambda}(X),
\]
the fuzzy Scott frame of \(X\), now regarded as a fuzzy domain via the internal
hom
\[
    [X,\Lambda]_{\mathrm{ff}}.
\]

If
\[
    f:X\to Y
\]
is a fuzzy-domain morphism, then
\[
    f^\ast:
    \mathcal O_{\Lambda}(Y)\to\mathcal O_{\Lambda}(X)
\]
is given by
\[
    f^\ast(V)=Vf.
\]

\section{The internal linear duality}
\label{sec:internal-linear-duality}

\subsection{The enriched construction}
\label{subsec:internal-linear-duality-enriched}

\begin{definition}[Linear dual]
\label{def:fuzzy-domain-linear-dual}
For a fuzzy domain \(X\), define its \emph{linear dual} by
\[
    X^{\perp}
    =
    [X,\underline{\mathbb V}]_{\mathrm{ff}}.
\]
Thus
\[
    X^{\perp}
    =
    \mathcal O_{\mathbb V}(X).
\]
\end{definition}

The assignment
\[
    X\longmapsto X^{\perp}
\]
is contravariant:
\[
    (-)^{\perp}:
    (\mathbf{FDom}_{\mathbb V}^{+})^{\op}
    \to
    \mathbf{FDom}_{\mathbb V}^{+}.
\]

The closed structure gives, for fuzzy domains \(A\) and \(X\), a natural
bijection
\[
    \mathbf{FDom}_{\mathbb V}^{+}
    (A,X^{\perp})
    \cong
    \mathbf{FDom}_{\mathbb V}^{+}
    (A\boxtimes_{\mathrm{ff}}X,\underline{\mathbb V}).
\]
By symmetry of
\[
    \boxtimes_{\mathrm{ff}},
\]
this is naturally isomorphic to
\[
    \mathbf{FDom}_{\mathbb V}^{+}
    (X,A^{\perp}).
\]

\begin{theorem}[Linear duality adjunction]
\label{thm:linear-duality-adjunction}
There is a contravariant self-adjunction
\[
    (-)^{\perp,\op}
    \dashv
    (-)^{\perp}
    :
    \mathbf{FDom}_{\mathbb V}^{+}
    \rightleftarrows
    (\mathbf{FDom}_{\mathbb V}^{+})^{\op}.
\]
Equivalently, for fuzzy domains \(A\) and \(X\), there is a natural bijection
\[
    \mathbf{FDom}_{\mathbb V}^{+}
    (A,X^{\perp})
    \cong
    \mathbf{FDom}_{\mathbb V}^{+}
    (X,A^{\perp}).
\]
\end{theorem}

\begin{proof}
Using closedness,
\[
\begin{aligned}
    \mathbf{FDom}_{\mathbb V}^{+}(A,X^{\perp})
    &=
    \mathbf{FDom}_{\mathbb V}^{+}
    (A,[X,\underline{\mathbb V}]_{\mathrm{ff}})
    \\
    &\cong
    \mathbf{FDom}_{\mathbb V}^{+}
    (A\boxtimes_{\mathrm{ff}}X,\underline{\mathbb V}).
\end{aligned}
\]
By symmetry,
\[
    A\boxtimes_{\mathrm{ff}}X
    \cong
    X\boxtimes_{\mathrm{ff}}A.
\]
Applying closedness again gives
\[
\begin{aligned}
    \mathbf{FDom}_{\mathbb V}^{+}
    (X\boxtimes_{\mathrm{ff}}A,\underline{\mathbb V})
    &\cong
    \mathbf{FDom}_{\mathbb V}^{+}
    (X,[A,\underline{\mathbb V}]_{\mathrm{ff}})
    \\
    &=
    \mathbf{FDom}_{\mathbb V}^{+}
    (X,A^{\perp}).
\end{aligned}
\]
These bijections are natural in \(A\) and \(X\), giving the claimed
contravariant adjunction.
\end{proof}

The transposition associated to a morphism
\[
    h:A\to X^{\perp}
\]
is the composite
\[
    A\boxtimes_{\mathrm{ff}}X
    \xrightarrow{h\boxtimes 1_X}
    X^{\perp}\boxtimes_{\mathrm{ff}}X
    \xrightarrow{\operatorname{ev}}
    \underline{\mathbb V},
\]
then re-curried, using symmetry, as a morphism
\[
    X\to A^{\perp}.
\]

\section{Fuzzy distributions as Scott double duals}
\label{sec:fuzzy-distributions-scott-double-duals}

\subsection{The enriched construction}
\label{subsec:fuzzy-distributions-enriched}

\begin{definition}[Fuzzy distribution object]
\label{def:fuzzy-distribution-object}
For a fuzzy domain \(X\), define its \emph{fuzzy distribution object} by
\[
    \mathbb D_{\mathbb V}X
    =
    X^{\perp\perp}
    =
    [X^{\perp},\underline{\mathbb V}]_{\mathrm{ff}}.
\]
Equivalently,
\[
    \mathbb D_{\mathbb V}X
    =
    [\mathcal O_{\mathbb V}(X),\underline{\mathbb V}]_{\mathrm{ff}}.
\]
\end{definition}

An object
\[
    \mu\in\mathbb D_{\mathbb V}X
\]
is a fuzzy-domain morphism
\[
    \mu:\mathcal O_{\mathbb V}(X)\to\underline{\mathbb V}.
\]
Thus a distribution is a finite-flat-continuous functional on fuzzy Scott
observations.

\begin{definition}[Integration pairing]
\label{def:fuzzy-integration-pairing}
The \emph{integration pairing} is the evaluation morphism
\[
    \mathbb D_{\mathbb V}X
    \boxtimes_{\mathrm{ff}}
    \mathcal O_{\mathbb V}(X)
    \longrightarrow
    \underline{\mathbb V}.
\]
For
\[
    \mu\in\mathbb D_{\mathbb V}X
\]
and
\[
    U\in\mathcal O_{\mathbb V}(X),
\]
we write
\[
    \int_X U\,d\mu
    =
    \mu(U).
\]
\end{definition}

The pairing is the evaluation map for the internal hom
\[
    [\mathcal O_{\mathbb V}(X),\underline{\mathbb V}]_{\mathrm{ff}}.
\]

\subsection{The posetal reflection}
\label{subsec:fuzzy-distributions-posetal}

For a fuzzy domain \(X\) over \(\Lambda\),
\[
    \mathbb D_{\Lambda}X
    =
    [\mathcal O_{\Lambda}(X),\Lambda]_{\mathrm{ff}}.
\]
Thus a fuzzy distribution is a fuzzy-domain morphism
\[
    \mu:\mathcal O_{\Lambda}(X)\to\Lambda.
\]
Equivalently, it is a finite-flat-colimit-preserving functional on the fuzzy
Scott frame of \(X\).

The hom-value between distributions
\[
    \mu,\nu:\mathcal O_{\Lambda}(X)\to\Lambda
\]
is
\[
    \mathbb D_{\Lambda}X(\mu,\nu)
    =
    \bigmeet_{U\in\mathcal O_{\Lambda}(X)}
    \bigl(\mu(U)\multimap\nu(U)\bigr).
\]

Integration is evaluation:
\[
    \int_X U\,d\mu
    =
    \mu(U),
    \qquad
    U\in\mathcal O_{\Lambda}(X).
\]

\section{Dirac distributions}
\label{sec:dirac-distributions-fuzzy-domain}

\subsection{The enriched construction}
\label{subsec:dirac-enriched}

The evaluation morphism
\[
    X^{\perp}\boxtimes_{\mathrm{ff}}X
    \longrightarrow
    \underline{\mathbb V}
\]
has, by closedness and symmetry, a transpose
\[
    \delta_X:
    X\to X^{\perp\perp}.
\]

\begin{definition}[Dirac distribution]
\label{def:fuzzy-dirac-distribution}
The \emph{Dirac map} of a fuzzy domain \(X\) is the morphism
\[
    \delta_X:
    X\to\mathbb D_{\mathbb V}X
\]
obtained by transposing evaluation.  Elementwise,
\[
    \delta_x(U)=U(x)
\]
for
\[
    U\in\mathcal O_{\mathbb V}(X).
\]
\end{definition}

\begin{proposition}[Dirac distributions are well-defined]
\label{prop:dirac-distributions-well-defined}
The map
\[
    \delta_X:X\to\mathbb D_{\mathbb V}X
\]
is a fuzzy-domain morphism.
\end{proposition}

\begin{proof}
The evaluation morphism
\[
    X^{\perp}\boxtimes_{\mathrm{ff}}X
    \to
    \underline{\mathbb V}
\]
is a morphism in
\[
    \mathbf{FDom}_{\mathbb V}^{+}.
\]
By the closed structure, its transpose
\[
    X\to[X^{\perp},\underline{\mathbb V}]_{\mathrm{ff}}
\]
is also a fuzzy-domain morphism.  Since
\[
    [X^{\perp},\underline{\mathbb V}]_{\mathrm{ff}}
    =
    \mathbb D_{\mathbb V}X,
\]
this transpose is precisely \(\delta_X\).
\end{proof}

\begin{proposition}[Dirac integration]
\label{prop:dirac-integration-fuzzy-domain}
For every fuzzy Scott observation
\[
    U\in\mathcal O_{\mathbb V}(X),
\]
one has
\[
    \int_X U\,d\delta_x
    =
    U(x).
\]
\end{proposition}

\begin{proof}
By definition,
\[
    \int_X U\,d\delta_x
    =
    \delta_x(U)
    =
    U(x).
\]
\end{proof}

\subsection{The posetal reflection}
\label{subsec:dirac-posetal}

For a fuzzy domain \(X\) over \(\Lambda\), the Dirac distribution at
\(x\in X\) is
\[
    \delta_x:\mathcal O_{\Lambda}(X)\to\Lambda,
    \qquad
    \delta_x(U)=U(x).
\]
Since finite-flat weighted colimits in
\[
    \mathcal O_{\Lambda}(X)
    =
    [X,\Lambda]_{\mathrm{ff}}
\]
are computed pointwise, evaluation at \(x\) preserves them.  Hence
\[
    \delta_x:\mathcal O_{\Lambda}(X)\to\Lambda
\]
is a fuzzy-domain morphism.

\section{Pushforward of fuzzy distributions}
\label{sec:pushforward-fuzzy-distributions}

\subsection{The enriched construction}
\label{subsec:pushforward-enriched}

Let
\[
    f:X\to Y
\]
be a fuzzy-domain morphism.  Pullback of observations gives
\[
    f^{\ast}:
    Y^{\perp}\to X^{\perp}.
\]
Dualizing again gives
\[
    (f^{\ast})^{\perp}:
    X^{\perp\perp}\to Y^{\perp\perp}.
\]

\begin{definition}[Pushforward]
\label{def:fuzzy-pushforward}
The \emph{pushforward} of fuzzy distributions along
\[
    f:X\to Y
\]
is the fuzzy-domain morphism
\[
    f_{\#}:
    \mathbb D_{\mathbb V}X\to\mathbb D_{\mathbb V}Y
\]
defined by
\[
    f_{\#}
    =
    (f^{\ast})^{\perp}.
\]
Elementwise,
\[
    f_{\#}(\mu)(V)
    =
    \mu(Vf)
\]
for
\[
    V\in\mathcal O_{\mathbb V}(Y).
\]
\end{definition}

\begin{proposition}[Functoriality of pushforward]
\label{prop:pushforward-functoriality-fuzzy-domain}
The assignment
\[
    X\longmapsto\mathbb D_{\mathbb V}X,
    \qquad
    f\longmapsto f_{\#}
\]
defines a covariant endofunctor
\[
    \mathbb D_{\mathbb V}:
    \mathbf{FDom}_{\mathbb V}^{+}
    \to
    \mathbf{FDom}_{\mathbb V}^{+}.
\]
\end{proposition}

\begin{proof}
The observation assignment is contravariant:
\[
    f:X\to Y
    \quad\mapsto\quad
    f^{\ast}:Y^{\perp}\to X^{\perp}.
\]
Applying the same contravariant duality again gives a covariant assignment
\[
    (f^{\ast})^{\perp}:X^{\perp\perp}\to Y^{\perp\perp}.
\]
For composable maps
\[
    X\xrightarrow{f}Y\xrightarrow{g}Z,
\]
one has
\[
    (gf)^{\ast}=f^{\ast}g^{\ast}.
\]
Dualizing gives
\[
    (gf)_{\#}
    =
    g_{\#}f_{\#}.
\]
The identity law is immediate.
\end{proof}

\begin{proposition}[Change of variables]
\label{prop:fuzzy-change-of-variables}
For every fuzzy-domain morphism
\[
    f:X\to Y,
\]
every distribution
\[
    \mu\in\mathbb D_{\mathbb V}X,
\]
and every fuzzy Scott observation
\[
    V\in\mathcal O_{\mathbb V}(Y),
\]
one has
\[
    \int_Y V\,d(f_{\#}\mu)
    =
    \int_X Vf\,d\mu.
\]
\end{proposition}

\begin{proof}
By definition,
\[
    \int_Y V\,d(f_{\#}\mu)
    =
    f_{\#}\mu(V)
    =
    \mu(Vf)
    =
    \int_X Vf\,d\mu.
\]
\end{proof}

\subsection{The posetal reflection}
\label{subsec:pushforward-posetal}

For a fuzzy-domain morphism
\[
    f:X\to Y
\]
over \(\Lambda\), pushforward is
\[
    f_{\#}(\mu)(V)
    =
    \mu(Vf),
    \qquad
    V\in\mathcal O_{\Lambda}(Y).
\]
This is a fuzzy-domain morphism
\[
    \mathbb D_{\Lambda}X\to\mathbb D_{\Lambda}Y.
\]

\section{The fuzzy distribution monad}
\label{sec:fuzzy-distribution-monad}

\subsection{Expectation observations}
\label{subsec:expectation-observations}

Let
\[
    X
\]
be a fuzzy domain.  For every fuzzy Scott observation
\[
    U\in X^{\perp},
\]
there is an associated fuzzy Scott observation on the distribution object
\[
    \widehat U:\mathbb D_{\mathbb V}X\to\underline{\mathbb V}
\]
defined by
\[
    \widehat U(\mu)=\mu(U).
\]

\begin{definition}[Expectation observation map]
\label{def:expectation-observation-map}
The \emph{expectation observation map} is the fuzzy-domain morphism
\[
    \operatorname{Exp}_X:
    X^{\perp}\to(\mathbb D_{\mathbb V}X)^{\perp}
\]
defined by
\[
    \operatorname{Exp}_X(U)=\widehat U.
\]
\end{definition}

Equivalently,
\[
    \operatorname{Exp}_X
    =
    \delta_{X^{\perp}}.
\]
Indeed, since
\[
    \mathbb D_{\mathbb V}X
    =
    X^{\perp\perp},
\]
we have
\[
    (\mathbb D_{\mathbb V}X)^{\perp}
    =
    X^{\perp\perp\perp},
\]
and for
\[
    U\in X^{\perp},
    \qquad
    \mu\in X^{\perp\perp},
\]
one has
\[
    \delta_{X^{\perp}}(U)(\mu)
    =
    \mu(U)
    =
    \widehat U(\mu).
\]

\begin{proposition}[Expectation observations are fuzzy Scott opens]
\label{prop:expectation-observations-are-scott}
The assignment
\[
    U\longmapsto\widehat U
\]
defines a fuzzy-domain morphism
\[
    \operatorname{Exp}_X:
    X^{\perp}\to(\mathbb D_{\mathbb V}X)^{\perp}.
\]
\end{proposition}

\begin{proof}
The evaluation morphism for the internal hom
\[
    \mathbb D_{\mathbb V}X
    =
    [X^{\perp},\underline{\mathbb V}]_{\mathrm{ff}}
\]
is
\[
    \mathbb D_{\mathbb V}X
    \boxtimes_{\mathrm{ff}}
    X^{\perp}
    \to
    \underline{\mathbb V}.
\]
Using the symmetry of
\[
    \boxtimes_{\mathrm{ff}},
\]
we also have a morphism
\[
    X^{\perp}
    \boxtimes_{\mathrm{ff}}
    \mathbb D_{\mathbb V}X
    \to
    \underline{\mathbb V}.
\]
By closedness, this transposes to
\[
    X^{\perp}
    \to
    [\mathbb D_{\mathbb V}X,\underline{\mathbb V}]_{\mathrm{ff}}
    =
    (\mathbb D_{\mathbb V}X)^{\perp}.
\]
This transpose is exactly
\[
    U\mapsto\widehat U.
\]
\end{proof}

\subsection{Monad multiplication}
\label{subsec:monad-multiplication-fuzzy-distributions}

\begin{definition}[Flattening]
\label{def:fuzzy-distribution-flattening}
Define
\[
    m_X:
    \mathbb D_{\mathbb V}\mathbb D_{\mathbb V}X
    \to
    \mathbb D_{\mathbb V}X
\]
by dualizing the expectation observation map:
\[
    m_X
    =
    (\operatorname{Exp}_X)^{\perp}.
\]
Elementwise, for
\[
    \Xi\in\mathbb D_{\mathbb V}\mathbb D_{\mathbb V}X
\]
and
\[
    U\in\mathcal O_{\mathbb V}(X),
\]
\[
    m_X(\Xi)(U)
    =
    \Xi(\widehat U).
\]
\end{definition}

\begin{theorem}[Fuzzy distribution monad]
\label{thm:fuzzy-distribution-monad}
The endofunctor
\[
    \mathbb D_{\mathbb V}
    =
    (-)^{\perp\perp}
\]
together with the Dirac maps
\[
    \delta_X:X\to\mathbb D_{\mathbb V}X
\]
and flattening maps
\[
    m_X:\mathbb D_{\mathbb V}\mathbb D_{\mathbb V}X
    \to
    \mathbb D_{\mathbb V}X
\]
forms a monad on
\[
    \mathbf{FDom}_{\mathbb V}^{+}.
\]
\end{theorem}

\begin{proof}
The monad is induced by the contravariant self-adjunction
\[
    (-)^{\perp,\op}
    \dashv
    (-)^{\perp}.
\]
The unit of the monad is the unit of this adjunction,
\[
    \delta_X:X\to X^{\perp\perp},
\]
which is the transpose of evaluation
\[
    X^{\perp}\boxtimes_{\mathrm{ff}}X
    \to
    \underline{\mathbb V}.
\]
The multiplication is the corresponding adjunction multiplication
\[
    X^{\perp\perp\perp\perp}
    \to
    X^{\perp\perp}.
\]
Under the explicit description of
\[
    \mathbb D_{\mathbb V}X
    =
    [X^{\perp},\underline{\mathbb V}]_{\mathrm{ff}},
\]
this multiplication is exactly
\[
    m_X(\Xi)(U)=\Xi(\widehat U).
\]

The monad laws follow from the triangle identities of the adjunction.  We also
record the elementwise verification.  For the left unit law,
\[
\begin{aligned}
    m_X(\delta_{\mu})(U)
    &=
    \delta_{\mu}(\widehat U)\\
    &=
    \widehat U(\mu)\\
    &=
    \mu(U).
\end{aligned}
\]
For the right unit law,
\[
\begin{aligned}
    m_X((\delta_X)_{\#}\mu)(U)
    &=
    ((\delta_X)_{\#}\mu)(\widehat U)\\
    &=
    \mu(\widehat U\delta_X)\\
    &=
    \mu(U),
\end{aligned}
\]
because
\[
    \widehat U(\delta_x)
    =
    \delta_x(U)
    =
    U(x).
\]

For associativity, let
\[
    \Theta\in \mathbb D_{\mathbb V}^{3}X,
    \qquad
    U\in X^{\perp}.
\]
Both composites
\[
    m_X m_{\mathbb D_{\mathbb V}X},
    \qquad
    m_X \mathbb D_{\mathbb V}(m_X)
\]
evaluate against \(U\) as follows:
\[
\begin{aligned}
    m_X(m_{\mathbb D_{\mathbb V}X}(\Theta))(U)
    &=
    m_{\mathbb D_{\mathbb V}X}(\Theta)(\widehat U)\\
    &=
    \Theta(\widehat{\widehat U}),
\end{aligned}
\]
and
\[
\begin{aligned}
    m_X(\mathbb D_{\mathbb V}(m_X)(\Theta))(U)
    &=
    \mathbb D_{\mathbb V}(m_X)(\Theta)(\widehat U)\\
    &=
    \Theta(\widehat U\,m_X)\\
    &=
    \Theta(\widehat{\widehat U}).
\end{aligned}
\]
Here
\[
    \widehat{\widehat U}(\Xi)
    =
    \Xi(\widehat U).
\]
Thus the two associativity composites agree.
\end{proof}

\subsection{Expectation of expectations}
\label{subsec:expectation-of-expectations}

For
\[
    \Xi\in\mathbb D_{\mathbb V}\mathbb D_{\mathbb V}X
\]
and
\[
    U\in\mathcal O_{\mathbb V}(X),
\]
flattening satisfies
\[
    \int_X U\,d\,m_X(\Xi)
    =
    \int_{\mu:\mathbb D_{\mathbb V}X}
    \left(
        \int_X U\,d\mu
    \right)
    d\Xi.
\]
Both sides are equal to
\[
    \Xi(\widehat U).
\]

\section{Fuzzy kernels and Chapman--Kolmogorov composition}
\label{sec:fuzzy-kernels-chapman-kolmogorov}

\subsection{The enriched construction}
\label{subsec:kernels-enriched}

\begin{definition}[Fuzzy kernel]
\label{def:fuzzy-kernel}
A \emph{fuzzy kernel} from \(X\) to \(Y\) is a fuzzy-domain morphism
\[
    k:X\to\mathbb D_{\mathbb V}Y.
\]
For \(x\in X\), the corresponding distribution on \(Y\) is denoted
\[
    k_x.
\]
\end{definition}

\begin{definition}[Kleisli extension]
\label{def:fuzzy-kleisli-extension}
Let
\[
    k:X\to\mathbb D_{\mathbb V}Y
\]
be a fuzzy kernel.  Its Kleisli extension is
\[
    k^{\sharp}:
    \mathbb D_{\mathbb V}X\to\mathbb D_{\mathbb V}Y
\]
defined by
\[
    k^{\sharp}
    =
    m_Y\circ\mathbb D_{\mathbb V}k.
\]
Elementwise,
\[
    k^{\sharp}(\mu)(V)
    =
    \mu(x\mapsto k_x(V))
\]
for
\[
    V\in\mathcal O_{\mathbb V}(Y).
\]
\end{definition}

The observation
\[
    x\longmapsto k_x(V)
\]
is fuzzy Scott because it is the composite
\[
    X
    \xrightarrow{k}
    \mathbb D_{\mathbb V}Y
    \xrightarrow{\widehat V}
    \underline{\mathbb V}.
\]
Thus the elementwise formula is well-defined.

\begin{definition}[Kleisli composite]
\label{def:fuzzy-kleisli-composite}
For fuzzy kernels
\[
    k:X\to\mathbb D_{\mathbb V}Y,
    \qquad
    \ell:Y\to\mathbb D_{\mathbb V}Z,
\]
their Kleisli composite is
\[
    \ell\star k
    =
    \ell^{\sharp}k.
\]
Explicitly,
\[
    (\ell\star k)_x(W)
    =
    k_x(y\mapsto \ell_y(W))
\]
for every
\[
    W\in\mathcal O_{\mathbb V}(Z).
\]
\end{definition}

The observation
\[
    y\longmapsto \ell_y(W)
\]
is the composite
\[
    Y
    \xrightarrow{\ell}
    \mathbb D_{\mathbb V}Z
    \xrightarrow{\widehat W}
    \underline{\mathbb V}.
\]
Hence it lies in
\[
    \mathcal O_{\mathbb V}(Y).
\]

\begin{theorem}[Fuzzy Chapman--Kolmogorov formula]
\label{thm:fuzzy-chapman-kolmogorov}
For kernels
\[
    k:X\to\mathbb D_{\mathbb V}Y,
    \qquad
    \ell:Y\to\mathbb D_{\mathbb V}Z,
\]
and every fuzzy Scott observation
\[
    W\in\mathcal O_{\mathbb V}(Z),
\]
one has
\[
    \int_Z W\,d(\ell\star k)_x
    =
    \int_Y
    \left(
        \int_Z W\,d\ell_y
    \right)
    dk_x.
\]
\end{theorem}

\begin{proof}
By definition,
\[
\begin{aligned}
    \int_Z W\,d(\ell\star k)_x
    &=
    (\ell\star k)_x(W)\\
    &=
    k_x(y\mapsto\ell_y(W))\\
    &=
    \int_Y
    \left(
        \int_Z W\,d\ell_y
    \right)
    dk_x.
\end{aligned}
\]
\end{proof}

\begin{proposition}[Associativity of kernel composition]
\label{prop:fuzzy-kernel-composition-associative}
Kleisli composition of fuzzy kernels is associative, and the Dirac kernel
\[
    \delta_X:X\to\mathbb D_{\mathbb V}X
\]
is a two-sided identity.
\end{proposition}

\begin{proof}
This is associativity and unitality of Kleisli composition for the monad
\[
    \mathbb D_{\mathbb V}.
\]
Explicitly, for kernels
\[
    k:X\to\mathbb D_{\mathbb V}Y,
    \qquad
    \ell:Y\to\mathbb D_{\mathbb V}Z,
    \qquad
    r:Z\to\mathbb D_{\mathbb V}W,
\]
one has
\[
    r\star(\ell\star k)
    =
    (r\star\ell)\star k.
\]
After evaluation at
\[
    Q\in\mathcal O_{\mathbb V}(W),
\]
both sides are given by
\[
    k_x
    \bigl(
        y\mapsto
        \ell_y
        (z\mapsto r_z(Q))
    \bigr).
\]
\end{proof}

\section{Barycenter algebras}
\label{sec:fuzzy-barycenter-algebras}

\subsection{The enriched construction}
\label{subsec:barycenter-enriched}

\begin{definition}[Fuzzy barycenter algebra]
\label{def:fuzzy-barycenter-algebra}
A \emph{fuzzy barycenter algebra} is an algebra for the monad
\[
    \mathbb D_{\mathbb V}
\]
on
\[
    \mathbf{FDom}_{\mathbb V}^{+}.
\]
Thus it consists of a fuzzy domain \(A\) equipped with a fuzzy-domain morphism
\[
    b_A:\mathbb D_{\mathbb V}A\to A
\]
such that
\[
    b_A\delta_A=\id_A
\]
and
\[
    b_A\mathbb D_{\mathbb V}(b_A)=b_A m_A.
\]
\end{definition}

\begin{definition}[Integration into a barycenter algebra]
\label{def:fuzzy-barycentric-integration}
Let
\[
    (A,b_A)
\]
be a fuzzy barycenter algebra.  For a fuzzy-domain morphism
\[
    f:X\to A
\]
and a distribution
\[
    \mu\in\mathbb D_{\mathbb V}X,
\]
define the barycentric integral by
\[
    \int_X f\,d\mu
    =
    b_A(f_{\#}\mu).
\]
Equivalently, the integration morphism is the composite
\[
    \mathbb D_{\mathbb V}X
    \xrightarrow{f_{\#}}
    \mathbb D_{\mathbb V}A
    \xrightarrow{b_A}
    A.
\]
\end{definition}

\begin{proposition}[Dirac law for barycentric integration]
\label{prop:fuzzy-barycentric-dirac-law}
For every fuzzy-domain morphism
\[
    f:X\to A,
\]
one has
\[
    \int_X f\,d\delta_x
    =
    f(x).
\]
\end{proposition}

\begin{proof}
Since
\[
    f_{\#}\delta_x=\delta_{f(x)},
\]
the result follows from the algebra unit law
\[
    b_A\delta_A=\id_A.
\]
\end{proof}

\begin{theorem}[Barycentric chain rule]
\label{thm:fuzzy-barycentric-chain-rule}
Let
\[
    (A,b_A)
\]
be a fuzzy barycenter algebra.  Let
\[
    k:X\to\mathbb D_{\mathbb V}Y
\]
be a fuzzy kernel, and let
\[
    g:Y\to A
\]
be a fuzzy-domain morphism.  Define
\[
    h:X\to A
\]
by
\[
    h(x)=\int_Y g\,dk_x.
\]
Equivalently,
\[
    h
    =
    X
    \xrightarrow{k}
    \mathbb D_{\mathbb V}Y
    \xrightarrow{g_{\#}}
    \mathbb D_{\mathbb V}A
    \xrightarrow{b_A}
    A.
\]
Then, for every
\[
    \mu\in\mathbb D_{\mathbb V}X,
\]
one has
\[
    \int_Y g\,d(k^{\sharp}\mu)
    =
    \int_X h\,d\mu.
\]
\end{theorem}

\begin{proof}
By definition,
\[
    k^{\sharp}
    =
    m_Y\mathbb D_{\mathbb V}k.
\]
Therefore
\[
\begin{aligned}
    \int_Y g\,d(k^{\sharp}\mu)
    &=
    b_A\bigl(g_{\#}(m_Y(\mathbb D_{\mathbb V}k(\mu)))\bigr).
\end{aligned}
\]
Naturality of the monad multiplication gives
\[
    g_{\#}m_Y
    =
    m_A\mathbb D_{\mathbb V}(g_{\#}).
\]
Hence the last expression is
\[
    b_A m_A
    \bigl(
        \mathbb D_{\mathbb V}(g_{\#})
        \mathbb D_{\mathbb V}k(\mu)
    \bigr).
\]
Using the algebra law
\[
    b_A m_A
    =
    b_A\mathbb D_{\mathbb V}(b_A),
\]
this becomes
\[
    b_A
    \mathbb D_{\mathbb V}(b_A g_{\#}k)(\mu)
    =
    b_A h_{\#}(\mu)
    =
    \int_X h\,d\mu.
\]
\end{proof}

\section{Comparison with fuzzy ideals and fuzzy Scott opens}
\label{sec:comparison-fuzzy-ideals-scott-opens}

The fuzzy distribution monad of this chapter does not use the full presheaf
completion
\[
    [X^{\op},\underline{\mathbb V}]
\]
as a free module.  The presheaf object belongs to the fuzzy ideal and predicate
doctrine.  Inside it, the finite-flat weights determine fuzzy ideals:
\[
    \operatorname{Idl}_{\mathbb V}(X)
    \subseteq
    [X^{\op},\underline{\mathbb V}].
\]

A fuzzy domain is a strict algebra for the fuzzy ideal monad, and its
finite-flat colimits are the fuzzy ideal joins.  A fuzzy Scott open is not an
arbitrary upper predicate
\[
    X\to\underline{\mathbb V}.
\]
It is an upper predicate preserving these fuzzy ideal joins.

Thus the observation object in this chapter is
\[
    \mathcal O_{\mathbb V}(X)
    =
    [X,\underline{\mathbb V}]_{\mathrm{ff}},
\]
not the raw predicate object
\[
    [X,\underline{\mathbb V}].
\]

The comparison is:
\[
\begin{tikzcd}[column sep=huge,row sep=large]
    \mathcal O_{\mathbb V}(X)
    =
    [X,\underline{\mathbb V}]_{\mathrm{ff}}
    \arrow[r,hook]
    &
    [X,\underline{\mathbb V}]
    \\
    \mathbb D_{\mathbb V}X
    =
    [\mathcal O_{\mathbb V}(X),\underline{\mathbb V}]_{\mathrm{ff}}
    \arrow[u,phantom,"\text{dualizes Scott opens}" description]
    &
    [X^{\op},\underline{\mathbb V}]
    \arrow[u,phantom,"\text{raw predicate/ideal ambient object}" description]
\end{tikzcd}
\]
The upper-left object is the fuzzy Scott topology of \(X\).  The lower-left
object is the fuzzy distribution object.  The presheaf object on the right is
not the basis of the distribution construction.

\section{The posetal reflection}
\label{sec:posetal-reflection-fuzzy-distributions}

Let
\[
    \Lambda=(L,\leq,\meet,\join,\tensor,\multimap,\bot,\top)
\]
be a fuzzy truth-value object.  For a fuzzy domain \(X\) over \(\Lambda\), with
algebra map
\[
    \alpha_X:\operatorname{Idl}_{\Lambda}(X)\to X,
\]
the fuzzy Scott topology is
\[
    \mathcal O_{\Lambda}(X)
    =
    \operatorname{Scott}_{\Lambda}(X).
\]
Its elements are upper fuzzy predicates
\[
    U:X\to L
\]
such that
\[
    U(\alpha_X(I))
    =
    \bigjoin_{x\in X} I(x)\tensor U(x)
\]
for every fuzzy ideal
\[
    I:X^{\op}\to\Lambda.
\]

The fuzzy distribution object is
\[
    \mathbb D_{\Lambda}X
    =
    [\mathcal O_{\Lambda}(X),\Lambda]_{\mathrm{ff}}.
\]
Thus a fuzzy distribution is a finite-flat-continuous functional
\[
    \mu:\mathcal O_{\Lambda}(X)\to\Lambda.
\]

The basic elementwise formulas are:
\[
    \int_X U\,d\mu
    =
    \mu(U),
\]
\[
    \delta_x(U)=U(x),
\]
\[
    f_{\#}(\mu)(V)=\mu(Vf),
\]
and
\[
    m_X(\Xi)(U)
    =
    \Xi(\mu\mapsto\mu(U)).
\]

\section{Examples}
\label{sec:fuzzy-distribution-examples}

\subsection{The crisp separated case}
\label{subsec:crisp-separated-distribution-recovery}

For
\[
    \Lambda=2,
\]
separated fuzzy domains are dcpos, and fuzzy-domain morphisms are
Scott-continuous maps.

For a dcpo \(X\), the fuzzy Scott topology becomes the ordinary Scott topology:
\[
    \mathcal O_2(X)=\sigma(X),
\]
regarded as a dcpo ordered by inclusion.  The distribution object is
\[
    \mathbb D_2X
    =
    [\sigma(X),2]_{\mathrm{ff}}.
\]
Thus a \(2\)-valued fuzzy distribution is a Scott-continuous map
\[
    \sigma(X)\to 2.
\]
Equivalently, it is a Scott-open subset of the dcpo
\[
    \sigma(X).
\]
In other words, it is a Scott-open family of Scott opens of \(X\).

The Dirac distribution at \(x\in X\) is
\[
    \delta_x(U)=1
    \quad\Longleftrightarrow\quad
    x\in U
\]
for Scott opens
\[
    U\in\sigma(X).
\]

This is not the finite-probability distribution monad and not the Giry monad.
It is the Scott-topological double-dual monad
\[
    X\longmapsto[\sigma(X),2]_{\mathrm{ff}}
\]
internal to dcpos.

\subsection{G\"odel truth values}
\label{subsec:godel-distribution-example}

For the G\"odel base, a fuzzy Scott open is a monotone degree-valued predicate
\[
    U:X\to[0,1]
\]
preserving fuzzy ideal joins by
\[
    U(\alpha_X(I))
    =
    \bigjoin_{x\in X}\min(I(x),U(x)).
\]
A distribution is a finite-flat-continuous functional
\[
    \mu:\mathcal O_{\Lambda}(X)\to[0,1].
\]

\subsection{\L ukasiewicz truth values}
\label{subsec:lukasiewicz-distribution-example}

For the \L ukasiewicz base,
\[
    a\tensor b=\max(0,a+b-1).
\]
The fuzzy Scott condition becomes
\[
    U(\alpha_X(I))
    =
    \bigjoin_{x\in X}
    \max(0,I(x)+U(x)-1).
\]
Distributions are finite-flat-continuous functionals on these fuzzy Scott
opens.

\subsection{Product truth values}
\label{subsec:product-distribution-example}

For the product base,
\[
    a\tensor b=ab.
\]
The fuzzy Scott condition becomes
\[
    U(\alpha_X(I))
    =
    \bigjoin_{x\in X}
    I(x)U(x).
\]
Distributions are product-linear Scott-continuous functionals on the fuzzy
Scott topology.

\subsection{The Lawvere base}
\label{subsec:lawvere-distribution-example}

For the Lawvere base
\[
    [0,\infty]^{\op},
\]
the tensor is addition and joins in the quantale are infima in the usual order.
Thus the fuzzy Scott condition takes the form
\[
    U(\alpha_X(I))
    =
    \inf_{x\in X}
    \bigl(I(x)+U(x)\bigr),
\]
written in the usual numerical order.  Internally this is the
\(\Lambda\)-join formula
\[
    \bigjoin_x I(x)\tensor U(x),
\]
with
\[
    \bigjoin
\]
computed in the reversed order.  A distribution is a tropical-linear
Scott-continuous functional on the domain of Scott-continuous potentials.

\subsection{Non-chain truth values}
\label{subsec:non-chain-distribution-example}

For a non-chain truth-value object such as a downset lattice
\[
    L=\operatorname{Down}(P),
\]
fuzzy Scott opens are \(L\)-valued predicates preserving fuzzy ideal joins.
Distributions are finite-flat-continuous functionals
\[
    \mathcal O_{\Lambda}(X)\to\Lambda.
\]
Thus the construction applies equally to multi-aspect or resource-indexed
truth values.

\section{Summary}
\label{sec:summary-fuzzy-distributions}

This chapter constructed the fuzzy distribution monad using the internal
linear duality of fuzzy domains.

The construction is not based on the raw presheaf object
\[
    [X^{\op},\underline{\mathbb V}]
\]
and not on the raw observation object
\[
    [X,\underline{\mathbb V}].
\]
Instead, for a fuzzy domain \(X\), the observation object is the fuzzy Scott
topology
\[
    \mathcal O_{\mathbb V}(X)
    =
    [X,\underline{\mathbb V}]_{\mathrm{ff}}.
\]
Its objects are the fuzzy Scott opens of \(X\), equivalently the truth-valued
fuzzy-domain morphisms
\[
    X\to\underline{\mathbb V}.
\]

The linear dual of a fuzzy domain is
\[
    X^{\perp}
    =
    [X,\underline{\mathbb V}]_{\mathrm{ff}}.
\]
The symmetric monoidal closed structure of fuzzy domains gives a contravariant
self-adjunction
\[
    (-)^{\perp,\op}
    \dashv
    (-)^{\perp}.
\]
The associated double-dual monad is
\[
    \mathbb D_{\mathbb V}X
    =
    X^{\perp\perp}
    =
    [[X,\underline{\mathbb V}]_{\mathrm{ff}},
      \underline{\mathbb V}]_{\mathrm{ff}}.
\]

A fuzzy distribution on \(X\) is therefore a fuzzy-domain morphism
\[
    \mu:\mathcal O_{\mathbb V}(X)\to\underline{\mathbb V}.
\]
Integration is evaluation:
\[
    \int_X U\,d\mu=\mu(U).
\]
The Dirac distribution is
\[
    \delta_x(U)=U(x).
\]
Pushforward along
\[
    f:X\to Y
\]
is
\[
    f_{\#}(\mu)(V)=\mu(Vf).
\]
Monad multiplication is expectation of expectations:
\[
    m_X(\Xi)(U)
    =
    \Xi(\mu\mapsto\mu(U)).
\]

Fuzzy kernels
\[
    k:X\to\mathbb D_{\mathbb V}Y
\]
compose by Kleisli composition, and satisfy the fuzzy
Chapman--Kolmogorov formula
\[
    \int_Z W\,d(\ell\star k)_x
    =
    \int_Y
    \left(
        \int_Z W\,d\ell_y
    \right)
    dk_x.
\]

In the posetal reflection, for a fuzzy domain \(X\) over \(\Lambda\),
\[
    \mathcal O_{\Lambda}(X)
    =
    \operatorname{Scott}_{\Lambda}(X),
\]
and
\[
    \mathbb D_{\Lambda}X
    =
    [\operatorname{Scott}_{\Lambda}(X),\Lambda]_{\mathrm{ff}}.
\]
In the crisp separated case, this becomes
\[
    \mathbb D_2X
    =
    [\sigma(X),2]_{\mathrm{ff}},
\]
the Scott-continuous double dual of the Scott topology of the dcpo \(X\).